%% kapa: exam=1 round=32 year=2021 period=1 form=A subject=civil_law subject_ko=민법 seq=27 number=27
%% source: Q-Net 감정평가사 기출문제 (cst00309) artlSeq=5208397 file=extracted/1008/5208397/A형/2021년도 제32회 감정평가사 1차 1교시 A형.pdf
%% conversion: kordoc 4.24.51 -> tools/gen1.py (auto); review=proofread-sonnet; flags=none
\begin{question}{27}
부동산 취득시효에 관한 설명으로 옳지 않은 것은? (다툼이 있으면 판례에 따름)
\begin{choices}
  \item 무과실은 점유취득시효의 요건이 아니다.
  \item 성명불상자의 소유물도 시효취득의 대상이 된다.
  \item 점유취득시효가 완성된 후에는 취득시효 완성의 이익을 포기할 수 있다.
  \item 행정재산은 시효취득의 대상이 아니다.
  \item 압류는 점유취득시효의 중단사유이다.
\end{choices}
\end{question}
\answer{5}
